\subsection{Fixed-objective convergence and numerical corners}
\label{sec:convergence-round3}
The primary recovered field is retained unchanged. A separate convergence investigation freezes its 32-node midpoint forward model, material-correction coefficients, 690 fitted bins, 147 supported cells, profiled pseudo-Huber log objective, graph weight and bounds. The diagnostic endpoints below are supplemental optimisation trials, not replacement reconstructions or evidence of improved physical accuracy.

For a selected numerical branch gradient $g_{\rm br}$, report
$\|w-\Pi_{[0,0.9]}(w-g_{\rm br})\|_\infty$ on supported cells separately from any nonsmooth stationarity statement. At the frozen endpoint this quantity is $3.3118970\times10^{-3}$, above the original $10^{-6}$ threshold. In cell 375, finite left and right signed secants at $h=10^{-8}$ are $-3.31397\times10^{-3}$ and $+4.09841\times10^{-3}$. Both feasible perturbations increase the objective in this coordinate. The selected-branch derivative follows the left trend; opposite one-sided secants are consistent with a numerical corner, not a unique smooth derivative. A coordinate check cannot establish constrained stationarity of the joint field.

Instrumented propagation reproduces every traced production flux to relative tolerance $10^{-13}$. Among 16 affected rays, the $\pm10^{-8}$ tests change one interpolation event: bin 831, quadrature node 16, final air step 1177. The actually used Air table crosses its 30\,TeV knot in the predictor and final-energy stopping-power evaluations. Adaptive step counts, limiting branches and air step floors are unchanged at this scale. Using the profiled robust envelope score, this ray accounts for $99.9529\%$ of the finite objective secant jump. This localises the numerical mechanism at the tested scales; it is not a derivative-limit proof.

A broader diagnostic evaluates 43 feasible coordinate and joint directions at six scales, with 22 of 258 trials giving resolved objective decreases. A first bounded polisher accepts 35 decreases in 827 trials and 20 full Jacobian assemblies. Its minimum-norm finite-gradient subproblem exposes an important distinction: small changes in simplex weights need not imply an accurately solved auxiliary convex problem. A revised pass therefore uses exact pairwise simplex line minimisation with an explicit Frank--Wolfe gap. Nearby-gradient methods motivate this investigation \citep{Burke2005Nonsmooth,Curtis2012Nonsmooth}, but the finite-budget heuristic used here does not inherit their convergence guarantees.

The revised pass also conditionally minimises the unchanged quadratic graph prior over 53 supported cells absent from every retained fitted ray. All eight unseen components are anchored; the harmonic extension preserves fitted fluxes and profiled normalisation exactly, with prior-gradient residual $3.55\times10^{-15}$. This leaves 94 directly constrained variables. Every sampled and trial field applies the harmonic extension; the reduced derivative is explicitly $T^{\mathsf T}g$, including floating-point residual contributions. Two radii, $10^{-6}$ and $10^{-8}$, each collect 21 gradients: the base, 12 seeded random feasible neighbours and both signs of four dominant coordinates. Auxiliary six-gradient and enriched hulls are compared at the same centre. The enriched pass accepts two exact decreases in eight line trials and 45 full assemblies. All accepted line trials decrease the original objective beyond a measured repeatability screen; no smoothing or tolerance relabelling changes the objective.

The frozen, first-polish and revised-polish objectives are respectively
$100.59269073659428$, $100.59269066369612$ and $100.59269064139075$.
Their selected-branch projected-gradient norms are $0.0033118970$, $0.0035453502$ and $0.0036712221$ (Figure~\ref{fig:convergence-round3}). The total objective decrease is only $9.52035\times10^{-8}$, below $10^{-9}$ relative, and the selected-branch criterion remains unmet. Independent full forward/Jacobian evaluations reproduce both polished endpoints, and source/runtime/output hashes pass. Accurate solutions of the finite sampled-hull convex problems do not certify the Clarke hull or constrained generalised stationarity of this 94-dimensional reduced problem. Termination is the preregistered finite budget; these trials neither establish calibrated inverse uncertainty nor alter the primary reconstruction.
\begin{figure}[t]\centering
\includegraphics[width=\textwidth]{revision/convergence_round3_20261025/fig_convergence_round3.pdf}
\caption{Supplemental fixed-objective convergence diagnostics. (a) Signed finite secants in coordinate 375 at the frozen endpoint; negative left and positive right values mean both tested perturbations raise the objective. Only the four smallest steps are plotted; the full six-step ladder is retained in the evidence. (b) Small independently reproduced objective decreases and selected-branch projected-gradient norms at three endpoints. The $10^{-6}$ branch criterion remains unmet; finite-hull QP accuracy is separate from generalised stationarity. No improvement in physical reconstruction is inferred from these tiny objective changes.}
\label{fig:convergence-round3}\end{figure}
